\chapter{Model Limitations and Honest Caveats}\label{chap:appE_limitations}

A research model earns trust by being explicit about what it does \emph{not} do.
The honest caveats that appear in context throughout the manual are consolidated
here as a single back-matter reference. None of them is hidden; each is documented
at its point of use and pointed to below. Read together, they fix the scientific
tier of \model{}: an EMIC-class, idealized coupled GCM for teaching and for
process and sensitivity studies, not a CMIP-class production Earth System Model.

% ----------------------------------------------------------------------
\section{Radiation and climate sensitivity}\label{sec:appE_radiation}
% ----------------------------------------------------------------------

\model{} has no multi-band radiative-transfer scheme and \emph{no top-of-atmosphere
energy closure}. The CO\textsubscript{2} response is injected as a single
surface-heat channel, so there is no Gregory regression, no effective radiative
forcing, and no feedback parameter $\lambda$. The reported
$0.43\,\mathrm{K\,(W\,m^{-2})^{-1}}$ is therefore a transient, TCR-like
surface-forcing \emph{proxy} computed about a (cold-tropics) base state, not a
calibrated equilibrium climate sensitivity (ECS) or transient climate response
(TCR): the water-vapour, lapse-rate, cloud and albedo feedbacks that amplify the
response in a full GCM are simply absent. The driver itself prints
``TCR-like proxy, NOT equilibrium ECS''. See
Chapter~\ref{chap:ch14_forcing_response} (\S\ref{sec:f14_honesty}) for the full
discussion.

% ----------------------------------------------------------------------
\section{Resolution}\label{sec:appE_resolution}
% ----------------------------------------------------------------------

The model runs at T31 ($\sim 3.75\degr$), chosen so it can integrate hundreds of
years per GPU-day. At this resolution it cannot resolve SST fronts, mesoscale
ocean eddies, tropical cyclones, or sharp orographic precipitation. Mesoscale
eddies are carried only through the Gent--McWilliams parameterization
(Chapter~\ref{chap:ch05_ocean_mixing}). The ocean dynamics carry a related
caveat. A full spherical \emph{prognostic} ocean core --- an implicit-viscous
no-slip baroclinic momentum balance, a Munk-boundary-layer transport
streamfunction, and GM, mass-conserving by construction --- is now implemented
and wired opt-in (\code{run\_dino\_control --prognostic-spherical}; see
Chapter~\ref{chap:ch07_ocean_prognostic_core}). It is genuinely density-driven:
the overturning emerges from the momentum balance, not from a box closure. Its
realistic \emph{magnitude}, however, is still being established through a free
(q-flux) spin-up --- on the fixed observed WOA density the geostrophic transport
overshoots to $\sim 160\Sv$, because the observed density has sharper gradients
than a coarse viscous model can sustain. The \emph{production} default therefore
remains the thermohaline closure, which at $130$-yr equilibrium gives a realistic
AMOC of $19.35\Sv$ at $26.5\degr$N (below-$500$~m, per-column metric; within the
RAPID observed $\sim 15$--$20\Sv$ range), cross-referenced in
Chapters~\ref{chap:ch08_ocean_overturning} and \ref{chap:ch17_amoc_tipping}.

% ----------------------------------------------------------------------
\section{Atmosphere}\label{sec:appE_atmosphere}
% ----------------------------------------------------------------------

The atmosphere uses single-humidity-column physics and produces a weak ITCZ
(precipitation too dry and displaced into the subtropics). The active multi-level
\dino{} spectral dycore fixes the structural gap of the legacy single-level
atmosphere --- it generates baroclinic eddies, jets and synoptic systems --- but
surface wind and mean-sea-level-pressure skill remains limited because the smoothed
T31 equator-to-pole SST gradient is too weak to build the observed mid-latitude
surface westerlies from the resolved eddy momentum flux
(Chapter~\ref{chap:ch18_jets_stormtracks}).

% ----------------------------------------------------------------------
\section{Ocean}\label{sec:appE_ocean}
% ----------------------------------------------------------------------

The ocean uses a linear equation of state, and the vertical bolus velocity
$w^\star$ is not formed from the Gent--McWilliams closed-form expression but
recovered from continuity of the combined (resolved-plus-bolus) flow, so GM enters
as a velocity rather than as a per-tracer skew-flux
(Chapter~\ref{chap:ch05_ocean_mixing}). The control climate is maintained by a
q-flux / flux-corrected restoring rather than by a fully free, CMIP-style
pre-industrial control: this pins the global mean and frees the SST enough to give
an honest forced response, but it is not an unconstrained \code{piControl}.

% ----------------------------------------------------------------------
\section{AMOC}\label{sec:appE_amoc}
% ----------------------------------------------------------------------

The Atlantic overturning is an \emph{interim} density-scaled thermohaline closure,
not resolved circulation. As its source docstring states plainly, it adds a
depth-integral-zero meridional overturning velocity over the Atlantic whose
strength scales with the subpolar-minus-subtropical upper-ocean density contrast,
with a \emph{prescribed} vertical shape but a physical strength and sensitivity and
a tunable bistability. AMOC tipping therefore does \emph{not} emerge from resolved
prognostic momentum; it is engineered into a Stommel/Marotzke-style box closure to
unblock density-responsiveness and tipping experiments
(Chapter~\ref{chap:ch17_amoc_tipping}). The fully prognostic-momentum core
(JEBAR barotropic-vorticity solver plus a $\dt{\uvec}$ momentum core) is still
under construction.
\codref{chronos\_esm/ocean/overturning.py}{interim density-scaled THC closure (P3/S1)}

% ----------------------------------------------------------------------
\section{Earth-system completeness}\label{sec:appE_completeness}
% ----------------------------------------------------------------------

\model{} is a coupled physical climate model, not yet a full Earth System Model in
the CMIP sense. Dissolved inorganic carbon is advected as a \emph{passive} tracer:
there is no air--sea gas exchange and no interactive carbon cycle. Sea ice uses
simplified Semtner thermodynamics with no ice \emph{dynamics}. The land surface is
a bucket model. There is no aerosol, ozone, volcanic or land-use forcing, which
limits fidelity for historical and scenario integrations. Atmospheric chemistry,
dynamic vegetation and ice sheets are absent or stubbed
(Chapter~\ref{chap:ch20_cmip7}, \S\ref{sec:c20_sci}).

% ----------------------------------------------------------------------
\section{Numerics}\label{sec:appE_numerics}
% ----------------------------------------------------------------------

To keep the model cleanly differentiable end to end, the iterative solvers use a
\emph{fixed} iteration count rather than data-dependent convergence loops: the
conjugate-gradient and \code{lax.scan}-based solvers run a frozen number of
iterations instead of \code{while}-loop tolerance tests. This makes the
computational graph static and reverse-mode-differentiable, at the cost of treating
solver convergence as a tuned constant rather than a checked condition.

% ----------------------------------------------------------------------

\begin{remark}[These limits are documented, not hidden]
Every caveat above is stated at its point of use in the body of the manual and,
where it is a code-level choice, in the source docstrings themselves. They define
the model's scientific scope rather than disqualify it: \model{} is positioned in
the EMIC / idealized-GCM tier, where a fast, fully differentiable, honestly
bounded coupled model is precisely the right instrument for teaching and for
process- and sensitivity-study research --- not a substitute for a CMIP-class
production ESM.
\end{remark}
